\documentclass[aps,pra,reprint,groupedaddress,nofootinbib,longbibliography]{revtex4-2}

\usepackage{amsmath,amssymb,bm,booktabs,array}
\usepackage[colorlinks=true,citecolor=blue,urlcolor=blue,linkcolor=blue]{hyperref}
\hypersetup{
  pdftitle={Exact unrestricted Gaussian rank of two canonical matchgate magic states},
  pdfauthor={Xiaojian Zeng}
}

\newcommand{\ket}[1]{\lvert #1\rangle}
\newcommand{\bra}[1]{\langle #1\rvert}
\newcommand{\C}{\mathbb{C}}
\newcommand{\rank}{\operatorname{rank}}
\newcommand{\Span}{\operatorname{span}}
\newcommand{\G}{\mathcal{G}}

\begin{document}

\title{Exact unrestricted Gaussian rank of two canonical matchgate magic states}

\author{Xiaojian Zeng}
\email{xiaojianzeng@foxmail.com}
\affiliation{PBC School of Finance, Tsinghua University, Beijing 100190, China}

\date{\today}

\begin{abstract}
Gaussian rank is the smallest number of pure fermionic Gaussian states whose
complex linear combination equals a target state, and it controls the branch
count in rank-based simulations of matchgate circuits with non-Gaussian inputs.
For the canonical four-mode state
$\ket{M}=(\ket{0000}+\ket{1111})/\sqrt{2}$, previous work established rank four
for $\ket{M}^{\otimes 2}$ only when every summand obeys two additional blockwise
symmetries, and conjectured the same value without those restrictions.  We prove
the unrestricted statement over $\C$.  After a parity reduction, a nonzero
Chevalley pairing places any hypothetical three-Gaussian decomposition, up to
$\operatorname{Spin}(16,\C)$, in one of 15 normal forms
$a\,1+b\,\Omega+c\,g_{q,i}$.  A two-form Clifford flattening has exact rank 78
on the target.  Exact characteristic-zero certificates put the 15 normal forms
on the two sides of a strict rank gap: eleven have rank at most 74 and four have
rank at least 83.  The certificate is checked by two separately implemented
exterior-algebra calculations.  Together with the four computational-basis
Gaussian summands, this gives $\chi_G(\ket{M}^{\otimes2})=4$ with no reality,
particle-number, copy-exchange, block-symmetry, or Pfaffian-chart assumption.
The result concerns exact rank for this finite eight-mode state; it does not
determine border rank, approximate rank, or many-copy asymptotics.
\end{abstract}

\maketitle

\section{Introduction}

Matchgate circuits and fermionic linear optics form a nontrivial class of quantum
processes admitting efficient classical descriptions
\cite{Valiant2002,TerhalDiVincenzo2002,Bravyi2005}.  Supplementing them with
non-Gaussian resources changes their computational power
\cite{BrodGalvao2011,Hebenstreit2019}.  One way to simulate such a supplemented
circuit is to decompose its input into pure fermionic Gaussian states and
propagate each Gaussian branch.  The minimum exact number of branches is the
Gaussian rank.

Cudby and Strelchuk introduced and studied this rank-based viewpoint for
matchgate magic states \cite{CudbyStrelchuk2025}.  For their canonical
four-mode state $\ket M$, they displayed the immediate four-term Gaussian
decomposition of two copies.  They proved the matching lower bound when each
Gaussian summand is constrained to the $+1$ eigenspaces of two blockwise
$ZZ$ operators.  Their numerical search found no unrestricted decomposition
with three terms, leading to the conjecture
\begin{equation}
  \chi_G(\ket M^{\otimes2})=4 .
  \label{eq:conjecture}
\end{equation}
The symmetry-restricted lower bound does not decide Eq.~\eqref{eq:conjecture},
because individual terms in a decomposition need not inherit a symmetry of the
sum.

Here we resolve Eq.~\eqref{eq:conjecture} without restrictions on the Gaussian
summands.  The proof combines a complete pure-spinor normal-form reduction with
an exact Clifford-flattening certificate.  Secant varieties of spinor varieties
have been studied at the level of dimensions and defectivity
\cite{Angelini2011}; those global dimensions do not by themselves decide whether
this particular target lies on the affine third secant.  The target-specific
rank gap below does.

The argument is organized as follows.  Section~\ref{sec:model} fixes the
fermionic conventions and the exact-rank question.  Section~\ref{sec:reduction}
reduces every possible three-term decomposition to 15 normal forms.  Sections
\ref{sec:flattening} and \ref{sec:certificate} construct and certify the rank
gap.  Section~\ref{sec:verification} records the computer-assisted component
and error control.  We close with the operational meaning and explicit scope
limitations.

\section{Eight-mode target and Gaussian rank}
\label{sec:model}

Let $\mathcal H_8=(\C^2)^{\otimes8}$ be the Fock space of eight ordered
fermionic modes.  We use the Jordan--Wigner Majoranas
\begin{align}
 c_{2j-1}&=Z_1\cdots Z_{j-1}X_j, &
 c_{2j}&=Z_1\cdots Z_{j-1}Y_j,
 \label{eq:majoranas}
\end{align}
for $j=1,\ldots,8$, so that $\{c_r,c_s\}=2\delta_{rs}I$.  Define
\begin{equation}
 \Lambda_8=\sum_{r=1}^{16}c_r\otimes c_r .
\end{equation}
In the intrinsic criterion used in Ref.~\cite{CudbyStrelchuk2025}, a normalized
pure state is fermionic Gaussian precisely when
\begin{equation}
 \Lambda_8(\ket g\otimes\ket g)=0 .
 \label{eq:gaussian-criterion}
\end{equation}
This definition allows either parity and does not assume particle-number
conservation.

For a nonzero state $\ket\phi$, its exact Gaussian rank over $\C$ is
\begin{equation}
 \begin{split}
 \chi_G(\phi)=\min\biggl\{r:\ &\ket\phi=\sum_{k=1}^r\alpha_k\ket{g_k},\\
 &\alpha_k\in\C,\quad \ket{g_k}\in\G_8\biggr\}.
 \end{split}
 \label{eq:rank-def}
\end{equation}
Normalization of each $\ket{g_k}$ causes no loss, since its scale can be
absorbed into $\alpha_k$.

The four-mode and eight-mode states of interest are
\begin{align}
 \ket M&=\frac{\ket{0000}+\ket{1111}}{\sqrt2},\\
 \ket\Psi=\ket M^{\otimes2}
 &=\frac12(\ket0+\ket{15}+\ket{240}+\ket{255}),
 \label{eq:target}
\end{align}
where the decimal labels denote the corresponding eight-bit Fock strings.
Every computational-basis vector is a pure Gaussian state, and direct
substitution in Eq.~\eqref{eq:gaussian-criterion} gives zero.  Equation
\eqref{eq:target} therefore proves $\chi_G(\Psi)\le4$.  We work below with the
unnormalized vector
\begin{equation}
 x=2\ket\Psi=\ket0+\ket{15}+\ket{240}+\ket{255},
 \label{eq:x}
\end{equation}
which has the same rank.

\section{Complete reduction of three Gaussian summands}
\label{sec:reduction}

Set $E=\C^8$ and use the exterior model
$S^+=\Lambda^{\mathrm{even}}E$ for the positive half-spin representation.
Nonzero even Gaussian vectors are exactly the affine cone over the even
pure-spinor variety \cite{Chevalley1954}.  Creation and annihilation operators
act as exterior multiplication $\epsilon_j$ and contraction $\iota_j$.

\subsection{Parity reduction}

Pure fermionic Gaussian vectors have definite parity.  If the even nonzero
vector $x$ were a sum of at most three definite-parity Gaussian vectors, its odd
projection would have to vanish.  One odd term must therefore have zero
coefficient; two odd terms must be proportional and cancel; and three odd terms
would leave no nonzero even part.  In the two-odd-term case the remaining even
part would be a single pure spinor.  More generally, the flattening introduced
in Sec.~\ref{sec:flattening} has rank 29 on each pure spinor and rank 78 on $x$,
so neither one nor two pure spinors can sum to $x$.  It is therefore sufficient
to exclude sums of three even pure spinors.

\subsection{A transverse pair}

Let $\Omega=e_1\wedge\cdots\wedge e_8$ and define the invariant Chevalley
pairing by
\begin{equation}
 B(s,t)=[\tau(s)\wedge t]_{\Lambda^8E},
 \label{eq:chevalley}
\end{equation}
where $\tau$ reverses wedge order.  A pure spinor pairs to zero with itself.
Direct evaluation of the four components in Eq.~\eqref{eq:x} gives
\begin{equation}
 B(x,x)=4\ne0 .
 \label{eq:target-pairing}
\end{equation}
Consequently, in any equality $x=s_1+s_2+s_3$ of pure spinors, at least one
pair $s_j,s_k$ has nonzero pairing.

For pure spinors, nonzero Chevalley pairing is equivalent to transversality of
their maximal-isotropic annihilators.  The complex spin group acts transitively
on ordered transverse pairs, taking their projective lines to those of $1$ and
$\Omega$.  The common projective stabilizer is the lifted $GL(E)$.  Hence a
hypothetical three-term decomposition can be transformed into
\begin{equation}
 u=a\,1+b\,\Omega+c\,g,
 \qquad a,b,c\in\C^*,
 \label{eq:three-normal}
\end{equation}
with $g$ an even pure spinor.  The coefficients are nonzero: if any vanished,
the flattening subadditivity bound would give rank at most $2\times29=58$, in
conflict with the target rank 78.

\subsection{Fifteen GL(8) orbits}

For completeness, the required pure-spinor normal form can be obtained directly
from its annihilator equations.  If $q$ is the lowest nonzero exterior degree
of $g$, its degree-$q$ component is decomposable.  A change of basis makes it
$e_1\wedge\cdots\wedge e_q$, and purity recursively gives
\begin{equation}
 g=e_1\wedge\cdots\wedge e_q\wedge\exp(\omega),
 \qquad
 \omega\in\Lambda^2\Span(e_{q+1},\ldots,e_8).
 \label{eq:pure-normal-general}
\end{equation}
Because $g\in S^+$, $q\in\{0,2,4,6,8\}$.  Complex skew congruence reduces
$\omega$ to a Darboux form of rank $2i$.  Absorbing the projective scale into
$c$ leaves
\begin{equation}
 g_{q,i}=e_1\cdots e_q
 \prod_{j=1}^{i}\left(1+e_{q+2j-1}e_{q+2j}\right),
 \quad 0\le i\le\frac{8-q}{2}.
 \label{eq:gqi}
\end{equation}
Thus $5+4+3+2+1=15$ cases exhaust all even pure spinors, including those with
zero vacuum amplitude.  No real structure or target symmetry has entered this
reduction.

\section{Two-form Clifford flattening}
\label{sec:flattening}

Let $V=E\oplus E^*$ be the 16-dimensional quadratic space and let
$\Gamma_1,\ldots,\Gamma_{16}$ be split Clifford generators.  For $u\in S^+$,
define
\begin{align}
 F_u &:S^+\longrightarrow\Lambda^2V^*,\\
 [F_u(y)]_{rs}&=B(u,\Gamma_r\Gamma_s y),\qquad r<s.
 \label{eq:flattening}
\end{align}
This is a $120\times128$ matrix in the natural exterior bases.  Invariance of
$B$ and covariance of Clifford multiplication imply
\begin{equation}
 \rank F_{h\cdot u}=\rank F_u,
 \qquad h\in\operatorname{Spin}(16,\C).
 \label{eq:equivariance}
\end{equation}
The physical Majoranas in Eq.~\eqref{eq:majoranas} correspond to
$\epsilon_j+\iota_j$ and $i(\epsilon_j-\iota_j)$.  Relative to the split
generators they only multiply rows of $F_u$ by nonzero powers of $i$, so the
rank is unchanged.

For the vacuum, only the degree-eight and degree-six columns supply complementary
top forms in Eq.~\eqref{eq:flattening}.  Direct basis action gives
\begin{equation}
 \rank F_1=1+\binom82=29.
 \label{eq:pure-rank}
\end{equation}
Equivariance then gives rank 29 for every nonzero even pure spinor.  This also
justifies the rank-at-most-58 assertions used above.

The flattening is used here as an orbit-separating covariant
\cite{Landsberg2012}, not as the usual
subadditive lower bound alone.  Indeed, the generic bound from three pure terms
would only be $3\times29=87$, whereas the target rank is 78.  The decisive fact
is that the complete normal forms in Eq.~\eqref{eq:gqi} avoid rank 78.

\section{Exact characteristic-zero rank-gap certificate}
\label{sec:certificate}

Rows are ordered lexicographically by the 120 pairs $r<s$, and columns by the
128 even occupation masks.  All entries of $F_x$ are integers.  The certificate
exhibits a $78\times78$ minor with
\begin{equation}
 \det=2361183241434822606848=2^{71},
\end{equation}
together with 50 linearly independent right-null vectors.  These give matching
lower and upper bounds,
\begin{equation}
 \rank F_x=78.
 \label{eq:target-rank}
\end{equation}

For each normal form
\begin{equation}
 u_{q,i}(a,b,c)=a\,1+b\,\Omega+c\,g_{q,i},
 \label{eq:uqiform}
\end{equation}
the same matrix is treated over the rational-function field $\mathbb Q(a,b,c)$.
Table~\ref{tab:gap} summarizes the certified alternatives.

\begin{table*}[t]
\caption{Exact bounds for all 15 normal forms with $abc\ne0$.  An upper bound
over $\mathbb Q(a,b,c)$ holds after every complex specialization because all
larger minors vanish identically.  Each lower bound is supplied by mutually
row- and column-disjoint nonzero minors.}
\label{tab:gap}
\begin{ruledtabular}
\begin{tabular}{c c c c}
Type & $(q,i)$ & Certified statement & Separation from $\rank F_x=78$\\
\hline
Low & $(0,0),(0,1),(2,0),(6,0),(6,1),(8,0)$ & $\rank F_{u_{q,i}}\le57$ & at most $74$\\
Low & $(0,2),(2,1),(4,0),(4,1),(4,2)$ & $\rank F_{u_{q,i}}\le74$ & at most $74$\\
High & $(0,3),(2,3)$ & $\rank F_{u_{q,i}}\ge84$ & at least $83$\\
High & $(0,4),(2,2)$ & $\rank F_{u_{q,i}}\ge83$ & at least $83$\\
\end{tabular}
\end{ruledtabular}
\end{table*}

For the eleven low cases, the certificate supplies 71 independent rational
right-null vectors when the upper bound is 57, and 54 when it is 74.  A possible
denominator zero creates no exceptional loophole: the nullspace identities over
$\mathbb Q(a,b,c)$ imply that every $58$- or $75$-minor, respectively, is the
zero polynomial.  Matrix rank cannot increase under specialization.

For the four high cases, the nonzero pattern separates into seven bipartite
row-column support components.  Selecting a minor in each component gives
disjoint rows and columns, so the minor sizes add.  The bounds are 84 for
$(0,3)$, 83 for $(0,4)$, 83 for $(2,2)$, and 84 for $(2,3)$.  All selected
determinants are nonzero monomials in $a,b,c$ except for the first component of
$(0,4)$.  There the certificate provides two alternative $11\times11$ minors,
\begin{align}
 -128a^2b^3c^5(a-b),\qquad
  64a^2b^3c^5(a+b).
 \label{eq:exception-minors}
\end{align}
For $abc\ne0$ in characteristic zero, $a-b$ and $a+b$ cannot vanish
simultaneously.  Thus the bound 83 holds on both exceptional loci as well as
generically.

We may now state the result.

\textit{Theorem.}  For $\ket M=(\ket{0000}+\ket{1111})/\sqrt2$, the exact
unrestricted fermionic Gaussian rank over $\C$ is
\begin{equation}
 \boxed{\chi_G(\ket M^{\otimes2})=4}.
\end{equation}

\textit{Proof.}  The four Fock-basis components in Eq.~\eqref{eq:target} are
Gaussian and give rank at most four.  Suppose the rank were at most three.
Parity reduces the putative decomposition to three even pure spinors.  Equations
\eqref{eq:target-pairing}--\eqref{eq:gqi} transform it, without changing
flattening rank, to one of the 15 matrices in Eq.~\eqref{eq:uqiform}.  Equation
\eqref{eq:target-rank} requires rank 78.  Table~\ref{tab:gap} and the exceptional
certificates \eqref{eq:exception-minors} instead give rank at most 74 or at
least 83 in every case, a contradiction.  Therefore the rank is at least four,
which meets the upper bound. \hfill$\square$

\section{Computer-assisted verification and error control}
\label{sec:verification}

The computer-assisted part is a finite exact certificate, not a numerical scan.
The generator constructs $F_u$ from bit-mask wedge and contraction operations,
computes the target pivot minor and nullspace, obtains rational-function
nullspaces for all low strata, and verifies every indexed high-stratum minor
before emitting a machine-readable record.  A second implementation reconstructs
the matrices through separately written wedge, contraction, and top-wedge
reversal operations.  It reads only the claimed certificate indices and vectors,
then recomputes every dot product, determinant, support component, and rank.
Both implementations use exact integer and rational-function arithmetic in
SymPy \cite{SymPy2017}.

OpenAI Codex workers assisted with certificate design and code generation.  The
registered model identifier was \texttt{gpt-5.6-sol}, with reasoning effort
\texttt{ultra}.  They were directed to reproduce the baseline, restrict the claim
to exact eight-mode rank, and require a separate replay.  Conclusions rest on the
retained exact identities, not model confidence; human inspection remains required.

The replay additionally substitutes directly on $a=b$ and $a=-b$ for the
$(0,4)$ stratum and obtains exact ranks 83 and 84, respectively.  It checks all
15 strata and reconstructs the target from the two four-mode factor supports.
The submitted certificate and scripts provide the row/column indices and sparse
null vectors; thus no trust in a reported rank routine is needed to check the
central identities.

There is no floating-point tolerance, statistical sampling, truncation, fitted
degree bound, finite-field lifting, or perturbative remainder.  A prior
finite-field scan was used only to discover the rank pattern and is not part of
the proof.  The analytic inputs that remain for human inspection are the parity,
transverse-pair, and pure-spinor normal-form reductions; the exact replay checks
the finite matrices conditional on those reductions.

\section{Physical meaning and limitations}
\label{sec:meaning}

The theorem closes the gap between the known four-branch decomposition and the
previous lower bound restricted by two blockwise $ZZ$ symmetries
\cite{CudbyStrelchuk2025}.  It shows that symmetry breaking, arbitrary complex
phases, non-number-conserving Gaussian states, and zero-vacuum-amplitude charts
cannot reduce the exact branch count below four for this two-copy input.  In a
simulation that propagates an exact Gaussian decomposition branch by branch, the
input therefore requires exactly four Gaussian branches.

This operational statement is deliberately finite.  Exact rank is not border
rank: the proof neither excludes a sequence of three-term states converging to
$\ket\Psi$ nor supplies an error-versus-cost bound for approximate simulation.
It also does not determine $\chi_G(\ket M^{\otimes k})$ for $k\ge3$, a
regularized rank, or an asymptotic runtime exponent.  The result is a property of
the specified eight-mode state and Gaussian resource set, not evidence for a new
particle, phase, or experimental effect.

\section{Conclusion}

An unrestricted rank-three decomposition of two canonical matchgate magic states
would, after an intrinsic pure-spinor reduction, have to occupy one of only 15
$GL(8)$ normal forms.  The two-form Clifford flattening places the target at rank
78 and every such form outside the interval from 75 through 82.  This exact rank
gap proves Gaussian rank four over $\C$ and removes all auxiliary symmetries from
the earlier two-copy lower bound.  The same strategy---orbit reduction followed
by a covariant rank gap---may be useful for other target-specific secant-membership
questions, although no many-copy claim follows from the present calculation.

\begin{acknowledgments}
In addition to the research uses disclosed in Sec.~\ref{sec:verification},
OpenAI Codex assisted with literature synthesis and drafting.  The tool is not an
author, and this draft requires the named author's independent review and
refinement before submission.
\end{acknowledgments}

\section*{Data availability}

The exact certificate and the two Python programs needed to regenerate and replay
the computer-assisted statements accompany this manuscript.  No experimental or
personal data were used.

\bibliography{references}

\end{document}
